\chapter{Effectful Cybernetic Mealy Systems}
\label{chap:effectful-cybernetic-mealy-systems}

\section{Orientation}
\label{sec:effectful-cybernetic-orientation}

The preceding chapters developed Mealy morphisms as typed stateful open systems.  A Mealy morphism
\[
\Phi:A\rightsquigarrow B
\]
has an input interface \(A\), an output interface \(B\), an internal state carrier, a transition
structure, and an output-comparison structure.  It is therefore a typed deterministic open process:
a current state and an admissible input event determine a next state together with a structured
output event.

This chapter adds effects to that process presentation and then uses the resulting effectful Mealy
systems to describe cybernetic feedback.  The guiding progression is
\[
\text{Mealy morphism}
\quad\leadsto\quad
\text{typed deterministic open process}
\quad\leadsto\quad
\text{effectful typed open process}
\quad\leadsto\quad
\text{feedback-closed cybernetic process}.
\]
Probability appears as one instance of the effect structure.  The first constructions do not depend
on probability.  They apply equally to nondeterminism, partiality, weighted computation,
probability, subprobability, and combinations of such effects whenever the corresponding response
effect has been constructed.

The central thesis of the chapter is:
\[
\boxed{
\text{An effectful cybernetic system is a feedback composite of effectful Mealy components.}
}
\]
In the probabilistic instance this becomes a typed categorical model of reinforcement-learning style
interaction:
\[
\boxed{
\text{environment feedback}
\quad+\quad
\text{stateful policy}
\quad+\quad
\text{optimizer/update dynamics}
\quad+\quad
\text{feedback closure}.
}
\]

There are four distinctions to keep separate.

First, a one-step transition-output assignment is not yet a Mealy morphism.  It becomes a Mealy
morphism only when it satisfies the identity and composite-input laws.

Second, an open Mealy process is not yet an autonomous process.  It has shape
\[
\text{state}+\text{admissible input}
\longmapsto
\text{next state}+\text{output}.
\]
Only after the input-output boundary is closed by a policy, controller, or environment does one
obtain an autonomous process.

Third, output is not automatically reward.  A Mealy morphism produces structured output in the
output interface.  A reward process is obtained only when part of this output is interpreted as a
reward, cost, performance signal, or value certificate.

Fourth, feedback trace supplies wiring.  It does not by itself create optimization or learning.
Those are stateful components of the controller.  The trace hides the plant-controller interaction
interface; the controller state carries policy, model, optimizer, and updater information.

This chapter does not prove convergence, optimality, Bellman optimality, regret bounds, or
statistical consistency.  It constructs the typed categorical process in which such learning
dynamics live.

\section{Deterministic Mealy Morphisms as Typed Open Processes}
\label{sec:deterministic-mealy-process-presentation}

Let \(E\) be the ambient category used for internal categories and spans in the preceding chapters.
Let
\[
\Phi:A\rightsquigarrow B
\]
be a Mealy morphism with carrier span
\[
A_0 \xleftarrow{p} P \xrightarrow{q} B_0 .
\]
The object \(P\) is the state object.  The maps \(p\) and \(q\) record the input-side and
output-side types exposed by a state.

\begin{definition}[Admissible state-input object]
\label{def:admissible-state-input-object}
The admissible state-input object of the carrier span \(P\) is
\[
D_P
=
A_1\times_{t_A,A_0,p}P .
\]
When \(P\) is the carrier of \(\Phi\), we also write
\[
D_\Phi=D_P .
\]
Thus an element of \(D_\Phi\) is a pair
\[
(a,x)
\]
where
\[
a:u\to p(x)
\]
is an input arrow of \(A\) whose target matches the input-side exposure of the state \(x\).
\end{definition}

Write
\[
\pi_A:D_\Phi\to A_1,
\qquad
\pi_x:D_\Phi\to P
\]
for the projections.

\begin{definition}[Typed next-output object]
\label{def:typed-next-output-object}
First define
\[
E_P
=
D_P
\times_{s_A\pi_A,A_0,p}
P .
\]
Thus \(E_P\) is the object of triples
\[
(a,x,y)
\]
satisfying
\[
t_A(a)=p(x),
\qquad
p(y)=s_A(a).
\]
Let
\[
\pi_y:E_P\to P
\]
denote the projection to the third component.  The typed next-output object of the carrier span
\(P\) is
\[
N_P
=
E_P
\times_{(q\pi_y,q\pi_x),\,B_0\times B_0,\,(s_B,t_B)}
B_1 .
\]
Equivalently, \(N_P\) is the object of tuples
\[
(a,x,y,\beta)
\]
satisfying
\[
t_A(a)=p(x),
\qquad
p(y)=s_A(a),
\qquad
s_B(\beta)=q(y),
\qquad
t_B(\beta)=q(x).
\]
The projection
\[
\rho_P:N_P\to D_P
\]
sends
\[
(a,x,y,\beta)\longmapsto(a,x).
\]
When \(P\) is the carrier of \(\Phi\), we write
\[
N_\Phi=N_P .
\]
\end{definition}

The deterministic Mealy structure assigns to each admissible pair \((a,x)\) a next state
\[
\delta_\Phi(a,x)\in P
\]
and an output comparison
\[
\omega_\Phi(a,x):q(\delta_\Phi(a,x))\to q(x)
\]
in \(B\).  The next state is typed by the source of the input arrow:
\[
p(\delta_\Phi(a,x))=s_A(a).
\]
Thus the Mealy structure determines a section
\[
\chi_\Phi:D_\Phi\to N_\Phi
\]
of \(\rho_\Phi:N_\Phi\to D_\Phi\), given elementwise by
\[
\chi_\Phi(a,x)
=
\bigl(a,x,\delta_\Phi(a,x),\omega_\Phi(a,x)\bigr).
\]

\begin{proposition}[One-step process presentation]
\label{prop:one-step-process-presentation}
Every Mealy morphism
\[
\Phi:A\rightsquigarrow B
\]
determines a typed deterministic one-step open process
\[
P
\xleftarrow{\;\pi_x\;}
D_\Phi
\xrightarrow{\;\delta_\Phi\;}
P
\]
together with a structured output map
\[
\omega_\Phi:D_\Phi\to B_1.
\]
Equivalently, it determines a section
\[
\chi_\Phi:D_\Phi\to N_\Phi
\]
of the typed next-output object.
\end{proposition}

\begin{proof}
The Mealy transition and output-comparison data assign to each admissible pair
\[
(a,x)\in A_1\times_{t_A,A_0,p}P
\]
a next state \(\delta_\Phi(a,x)\) and an output arrow
\[
\omega_\Phi(a,x):q(\delta_\Phi(a,x))\to q(x).
\]
The typing equations are exactly the equations defining \(N_\Phi\).  Hence the assignment
\[
(a,x)\mapsto
\bigl(a,x,\delta_\Phi(a,x),\omega_\Phi(a,x)\bigr)
\]
is a section of \(N_\Phi\to D_\Phi\).
\end{proof}

The converse requires equations.  A section \(D_\Phi\to N_\Phi\) is merely one-step data unless it
is compatible with identities and composite input events.

\begin{proposition}[Mealy laws in process form]
\label{prop:mealy-laws-process-form}
Let
\[
A_0\xleftarrow{p}P\xrightarrow{q}B_0
\]
be a carrier span.  A section
\[
\chi:D_P\to N_P
\]
with components
\[
(a,x)\longmapsto(\delta(a,x),\omega(a,x))
\]
defines a Mealy morphism structure on the carrier span precisely when the following equations hold.

For every state \(x\),
\[
\delta(1_{p(x)},x)=x,
\qquad
\omega(1_{p(x)},x)=1_{q(x)}.
\]

For composable input arrows
\[
b:w\to u,
\qquad
a:u\to p(x),
\]
one has
\[
\delta(a\circ b,x)
=
\delta\bigl(b,\delta(a,x)\bigr),
\]
and
\[
\omega(a\circ b,x)
=
\omega(a,x)\circ
\omega\bigl(b,\delta(a,x)\bigr).
\]
\end{proposition}

\begin{proof}
The identity equations express the unit law for the input-processing monad
\[
R_A=-\circ A
\]
acting on the carrier.  The composite-input equations express the algebra multiplication law for
\(R_A\), with the output comparison composed in \(B\).  Conversely, these equations give precisely
the unit and multiplication compatibility required for an \(R_A\)-algebra in the Kleisli category
of the output-comparison monad
\[
T_B=B\circ - .
\]
Thus the section defines the original Mealy morphism data exactly when these laws hold.
\end{proof}

The process is open.  It does not define a map
\[
P\to P
\]
because no input arrow has been chosen.  It instead defines a deterministic controlled transition
whose admissible input object at \(x\) is the fibre
\[
\{a:u\to p(x)\}.
\]

\subsection{One-object specialization}
\label{subsec:one-object-deterministic-process-specialization}

When \(A\) and \(B\) are one-object categories, their arrows form monoids.  In the free-monoid
case
\[
A=\Sigma^\ast,
\qquad
B=\Gamma^\ast,
\]
a classical deterministic Mealy machine has state set \(S\), transition map
\[
\tau:S\times \Sigma\to S,
\]
and output map
\[
\lambda:S\times \Sigma\to \Gamma .
\]
The generator-level one-step process presentation is
\[
S\times \Sigma\to S\times \Gamma,
\qquad
(s,a)\mapsto(\tau(s,a),\lambda(s,a)).
\]
The full word-level Mealy morphism is obtained when these one-step maps extend coherently along
the free monoid, equivalently when the identity and composite-input equations above hold.

\section{Indexed Response Effects}
\label{sec:indexed-response-effects}

The deterministic one-step process chooses a single next-state/output pair.  An effectful process
replaces this choice by an effectful response.  The correct place to add the effect is not only the
next state and not only the output; it is the typed next-state/output object.

The effect must also be indexed.  The response object
\[
N_P\to D_P
\]
is a displayed family of possible responses over the displayed family of admissible state-inputs.
Moreover, the composite-input law first samples one response and then uses the sampled next state
to choose the base of the second response.  Thus ordinary Kleisli extension in a single global
category is not enough; one needs dependent Kleisli substitution over slices.

\begin{definition}[Response-effect context]
\label{def:response-effect-context}
A response-effect context on \(E\) consists of an indexed Kleisli doctrine
\[
X\longmapsto \mathsf M_X:E/X\to E/X
\]
with the following structure.

\begin{enumerate}
\item For each object \(X\in E\), \(\mathsf M_X\) is a monad on the slice category \(E/X\).

\item For each map \(f:X'\to X\), there is a reindexing comparison
\[
f^\ast\mathsf M_X
\Longrightarrow
\mathsf M_{X'}f^\ast
\]
compatible with the units and Kleisli extensions.

\item Dependent Kleisli substitution is available: if
\[
Y\to X
\]
is a displayed response family and
\[
Z\to Y
\]
is a further displayed response family depending on the first response, then sections
\[
r:X\to \mathsf M_X(Y),
\qquad
k:Y\to \mathsf M_Y(Z)
\]
compose to a section
\[
r\bind k:X\to \mathsf M_X(Z).
\]
\end{enumerate}

The object
\[
\mathsf M_X(Y)\to X
\]
is read as the family of effectful responses over the same base \(X\).
\end{definition}

Thus, if
\[
Y\to X
\]
is a displayed family of possible responses over inputs \(X\), then
\[
\mathsf M_X(Y)\to X
\]
is the displayed family of effectful responses over the same inputs.

When \(E\) has a terminal object \(1\), we write
\[
\mathsf M=\mathsf M_1
\]
for the global effect on \(E\cong E/1\).

\subsection{Basic examples}
\label{subsec:basic-response-effects}

In finite sets, if \(Y\to X\) has fibres \(Y_x\), then the finite-distribution response effect has
fibres
\[
\mathsf M_X(Y)_x
=
\mathcal D(Y_x).
\]
The finite powerset response effect has fibres
\[
\mathsf M_X(Y)_x
=
\mathcal P_{\mathrm f}(Y_x),
\]
the maybe response effect has fibres
\[
\mathsf M_X(Y)_x
=
Y_x+1,
\]
and a finite subprobability response effect has fibres given by finite subprobability distributions
on \(Y_x\).

The intended readings are:
\[
\begin{array}{c|l}
\text{effect} & \text{meaning of an effectful response}\\
\hline
\mathrm{Id} & \text{a single deterministic response}\\
\mathcal P_{\mathrm f} & \text{a finite set of possible responses}\\
(-)+1 & \text{a response or failure}\\
\mathcal D & \text{a finite probability distribution over responses}\\
\mathcal S & \text{a finite subprobability distribution, allowing divergence or nontermination}\\
\mathcal W & \text{a finite weighted family of responses over a chosen semiring}
\end{array}
\]

Probability is therefore not a change to the underlying Mealy shape.  It is one choice of response
effect.

\section{Effectful Mealy Processes and Effectful Mealy Morphisms}
\label{sec:effectful-mealy-morphisms}

Let \(A\) and \(B\) be internal categories and let
\[
A_0\xleftarrow{p}P\xrightarrow{q}B_0
\]
be a carrier span.

\begin{definition}[Effectful one-step Mealy process]
\label{def:effectful-one-step-mealy-process}
An \(\mathsf M\)-effectful one-step Mealy process from \(A\) to \(B\) on the carrier span \(P\)
is a section
\[
\kappa:
D_P
\longrightarrow
\mathsf M_{D_P}(N_P)
\]
of the effectful response family over \(D_P\).
\end{definition}

Elementwise, for
\[
(a,x)\in D_P,
\qquad
a:u\to p(x),
\]
the value
\[
\kappa(a,x)
\]
is an effectful response over the fibre
\[
\{(y,\beta)\mid p(y)=u,\;\beta:q(y)\to q(x)\}.
\]
Thus, in element notation,
\[
\boxed{
\kappa(a,x)
\in
\mathsf M
\{(y,\beta):p(y)=s_A(a),\;\beta:q(y)\to q(x)\}.
}
\]

A one-step effectful response is not yet an effectful Mealy morphism.  It must satisfy the
effectful versions of the identity and composite-input laws.

\begin{definition}[Composable state-input object]
\label{def:composable-state-input-object}
Let
\[
A_0\xleftarrow{p}P\xrightarrow{q}B_0
\]
be a carrier span.  The composable state-input object is
\[
C_P
=
\bigl(A_1\times_{s_A,A_0,t_A}A_1\bigr)
\times_{t_A\pi_1,A_0,p}
P .
\]
Thus an element of \(C_P\) is a triple
\[
(a,b,x)
\]
where
\[
b:w\to u,
\qquad
a:u\to p(x).
\]
The composite input is
\[
a\circ b:w\to p(x).
\]
\end{definition}

\begin{definition}[Effectful Mealy morphism]
\label{def:effectful-mealy-morphism}
An \(\mathsf M\)-effectful Mealy morphism
\[
\Phi:A\rightsquigarrow_{\mathsf M}B
\]
consists of a carrier span
\[
A_0\xleftarrow{p}P\xrightarrow{q}B_0
\]
and an effectful one-step response
\[
\kappa_\Phi:D_\Phi\to\mathsf M_{D_\Phi}(N_\Phi)
\]
satisfying the following laws.

For every state \(x\),
\[
\kappa_\Phi(1_{p(x)},x)
=
\eta\bigl(x,1_{q(x)}\bigr),
\]
where the right-hand side denotes the unit response at the point
\[
(1_{p(x)},x,x,1_{q(x)})\in N_\Phi .
\]

The composite-input law is an equality of sections over the composable state-input object
\[
C_\Phi=C_P .
\]
Elementwise, for composable input arrows
\[
b:w\to u,
\qquad
a:u\to p(x),
\]
one has
\[
\kappa_\Phi(a\circ b,x)
=
\kappa_\Phi(a,x)
\bind
\Bigl((y,\beta)\mapsto
\kappa_\Phi(b,y)
\bind
\bigl((z,\gamma)\mapsto
\eta(z,\beta\circ\gamma)\bigr)
\Bigr).
\]
The right-hand side is formed by dependent Kleisli substitution: after sampling
\[
(y,\beta):q(y)\to q(x)
\]
from the response to \((a,x)\), the second response is taken at the state \(y\) and input \(b\).
\end{definition}

Thus the deterministic principle
\[
\text{state plus typed input}
\longmapsto
\text{next state plus typed output}
\]
is replaced by
\[
\text{state plus typed input}
\longmapsto
\text{effectful next-state/output response}.
\]

\begin{theorem}[Unit embedding of deterministic Mealy morphisms]
\label{thm:deterministic-mealy-embedding-effectful}
Every deterministic Mealy morphism
\[
\Phi:A\rightsquigarrow B
\]
determines an \(\mathsf M\)-effectful Mealy morphism
\[
\eta\Phi:A\rightsquigarrow_{\mathsf M}B
\]
by
\[
\kappa_{\eta\Phi}(a,x)
=
\eta
\bigl(\delta_\Phi(a,x),\omega_\Phi(a,x)\bigr).
\]
\end{theorem}

\begin{proof}
The deterministic Mealy morphism gives a section
\[
\chi_\Phi:D_\Phi\to N_\Phi .
\]
Composing this section with the unit of the slice effect
\[
\eta:N_\Phi\to \mathsf M_{D_\Phi}(N_\Phi)
\]
gives the required effectful response map.  The deterministic identity law is sent to the
effectful identity law by the unit.  For a composite input \(a\circ b\), the deterministic
composite-input equation says that the unique response to \(a\circ b\) is obtained by first taking
the unique response to \(a\), then the unique response to \(b\) at the resulting state, and then
composing the two output comparisons.  Applying the unit of \(\mathsf M\) turns this deterministic
composition into the displayed dependent Kleisli composite.  Hence \(\eta\Phi\) is an effectful
Mealy morphism.
\end{proof}

\section{Effectful Output-Comparison Kleisli Structure}
\label{sec:effectful-output-comparison-kleisli-structure}

The deterministic output-comparison monad from the Mealy-morphism calculus is
\[
T_B=B\circ - .
\]
The input-processing monad is
\[
R_A=-\circ A .
\]
A deterministic Mealy morphism may be regarded as an \(R_A\)-algebra in the Kleisli category of
\(T_B\).

To add effects, one should not assert that an arbitrary monad composed with \(T_B\) is automatically
again a monad.  The required object is an effectful output-comparison Kleisli structure.

\begin{definition}[Effectful output-comparison Kleisli structure]
\label{def:effectful-output-comparison-kleisli-structure}
An effectful output-comparison Kleisli structure for \(B\) consists of the following data.

\begin{enumerate}
\item A monad
\[
T_B^{\mathsf M}
\]
on the relevant local span category
\[
\Span(E)(A_0,B_0).
\]

\item A comparison of monads
\[
T_B\longrightarrow T_B^{\mathsf M}
\]
embedding deterministic output-comparison cells as unit-valued effectful output-comparison cells.

\item A lift of the input-processing monad
\[
R_A=-\circ A
\]
to the Kleisli category
\[
\Kl(T_B^{\mathsf M}).
\]
\end{enumerate}
\end{definition}

On the process presentation, \(T_B^{\mathsf M}\) sends the deterministic next-output object
\[
N_\Phi\to D_\Phi
\]
to the effectful next-output object
\[
\mathsf M_{D_\Phi}(N_\Phi)\to D_\Phi .
\]
In bicategorical notation this is written suggestively as
\[
T_B^{\mathsf M}P
=
\mathsf M(B\circ P),
\]
where \(\mathsf M\) is understood as acting in the appropriate displayed, indexed, or hom-category.
The notation \(T_B^{\mathsf M}\) records the structure actually needed: an effectful version of
output comparison together with a Kleisli action of input processing.

\begin{proposition}[Kleisli-algebra form]
\label{prop:effectful-kleisli-algebra-form}
Given an effectful output-comparison Kleisli structure for \(B\), an
\(\mathsf M\)-effectful Mealy morphism
\[
\Phi:A\rightsquigarrow_{\mathsf M}B
\]
is equivalently an algebra for the lifted input-processing monad
\[
R_A
\]
in
\[
\Kl(T_B^{\mathsf M}).
\]
Equivalently, it is a Kleisli algebra structure
\[
\alpha_\Phi:
R_A P
\longrightarrow
T_B^{\mathsf M}P,
\]
or, informally,
\[
P\circ A
\longrightarrow
\mathsf M(B\circ P).
\]
\end{proposition}

\begin{proof}
The object \(R_A P=P\circ A\) is the typed state-input object.  The object
\(T_B^{\mathsf M}P\) is the effectful typed next-state/output-comparison object.  Thus a Kleisli
arrow
\[
R_A P\to T_B^{\mathsf M}P
\]
is exactly an effectful one-step response.  The algebra unit law is the identity-input law, and the
algebra multiplication law is the dependent Kleisli composite-input law.
\end{proof}

When the response effect comes from a monad that distributes over output comparison, the structure
\(T_B^{\mathsf M}\) may be constructed as a genuine composite
\[
\mathsf M\circ T_B .
\]
The definition above isolates the structure needed for the Mealy calculus itself.

\section{Effectful Mealy Simulations}
\label{sec:effectful-mealy-simulations}

The previous chapters treat not only Mealy morphisms but also maps between them.  In the Mealy
simulation convention used here, simulations are oppositely directed on state carriers.

\begin{definition}[Strict effectful Mealy simulation]
\label{def:strict-effectful-mealy-simulation}
Let
\[
\Phi,\Psi:A\rightsquigarrow_{\mathsf M}B
\]
be effectful Mealy morphisms with carrier spans
\[
A_0\xleftarrow{p_P}P\xrightarrow{q_P}B_0,
\qquad
A_0\xleftarrow{p_Q}Q\xrightarrow{q_Q}B_0.
\]
A strict effectful Mealy simulation
\[
\Theta:\Phi\Rightarrow\Psi
\]
is a map
\[
\theta:Q\to P
\]
over \(A_0\) and \(B_0\),
\[
p_P\theta=p_Q,
\qquad
q_P\theta=q_Q,
\]
such that the square of effectful responses commutes:
\[
\mathsf M(N_\theta)\circ \kappa_\Psi
=
\kappa_\Phi\circ D_\theta .
\]
Here
\[
D_\theta:D_\Psi\to D_\Phi
\]
is induced by \(\theta\), and
\[
N_\theta:N_\Psi\to N_\Phi
\]
sends
\[
(a,x,y,\beta)
\longmapsto
(a,\theta x,\theta y,\beta).
\]
\end{definition}

The map \(N_\theta\) is well typed because \(\theta\) lies over \(B_0\).  If
\[
\beta:q_Q(y)\to q_Q(x),
\]
then
\[
q_P(\theta y)=q_Q(y),
\qquad
q_P(\theta x)=q_Q(x),
\]
so the same output-comparison arrow \(\beta\) may be read as
\[
\beta:q_P(\theta y)\to q_P(\theta x).
\]

\begin{proposition}[Deterministic simulations embed strictly]
\label{prop:deterministic-simulations-embed}
A strict deterministic Mealy simulation
\[
\Theta:\Phi\Rightarrow\Psi
\]
with carrier component
\[
\theta:Q\to P
\]
embeds as a strict effectful Mealy simulation
\[
\eta\Theta:\eta\Phi\Rightarrow\eta\Psi .
\]
\end{proposition}

\begin{proof}
The deterministic simulation equation says that transporting the response of \(\Psi\) along
\(\theta\) agrees with the response of \(\Phi\) after transporting the current state along
\(\theta\).  Applying the unit of the response effect gives
\[
\mathsf M(N_\theta)\circ \kappa_{\eta\Psi}
=
\kappa_{\eta\Phi}\circ D_\theta .
\]
\end{proof}

More general output-lax simulations are obtained algebraically as Kleisli algebra morphisms for
\(R_A\) in
\[
\Kl(T_B^{\mathsf M}),
\]
whenever the effectful output-comparison Kleisli structure is available.  The strict simulations
above are the identity-output-comparison special case.

\section{Abstract Cybernetic Mealy Architectures}
\label{sec:abstract-cybernetic-mealy-architectures}

We now move from single open processes to feedback architectures.

A cybernetic system contains a regulated component and a regulating component.  We call the
regulated component the plant and the regulating component the adaptive controller.  The plant
consumes actions and produces observations and performance signals.  The controller consumes
observations and performance signals and produces actions, together with any visible certificates,
values, or traces.

Let
\[
A=\text{action interface},
\]
\[
O=\text{observation interface},
\]
\[
R=\text{reward, cost, or performance interface},
\]
and
\[
V=\text{visible value, certificate, or trace interface}.
\]

In the concrete adaptive-controller presentation below, we write
\[
V=W\times T,
\]
where \(W\) is the value or certificate object and \(T\) is the trace or auxiliary-output object.
Thus the controller's visible output interface \(V\) records both value/certificate data and
policy-trace data.

There are two useful operational time conventions.

In the previous-feedback convention, the controller has type
\[
\mathbf K_{\mathrm{op}}:O\times R\rightsquigarrow A\times V.
\]
The input at time \(t\) is read as
\[
(o_t,r_{t-1}).
\]
Thus the reward or performance input is the feedback signal available before the current action is
chosen.

In the decision-update convention used for the closed-loop construction below, the controller is
decomposed into a within-step policy/evaluator and an across-step updater.  The policy chooses
using the current observation, and the updater changes the controller state after the plant returns
the next state and performance signal.

\subsection{Operational cybernetic architecture}
\label{subsec:operational-cybernetic-architecture}

The compact operational plant and controller have types
\[
\mathbf P:A\rightsquigarrow O\times R,
\qquad
\mathbf K_{\mathrm{op}}:O\times R\rightsquigarrow A\times V.
\]
The operational feedback loop is the synchronous closed-loop process obtained by feeding the action
output of the controller into the plant and the observation/performance output of the plant back
into the controller.  In the decision-update convention, this feedback is implemented by the
policy--plant--updater composite constructed in
Section~\ref{sec:closed-loop-process-presentation-effectful}.

\subsection{Signed-interface cybernetic architecture}
\label{subsec:signed-interface-cybernetic-architecture}

In the signed-interface convention, an Int morphism
\[
(X^+,X^-)\to(Y^+,Y^-)
\]
is represented by a Mealy morphism
\[
X^+ + Y^-\rightsquigarrow Y^+ + X^- .
\]
Thus the plant
\[
\mathbf P:A\rightsquigarrow O\times R
\]
represents an Int morphism
\[
\mathbf P:(0,0)\to(O\times R,A).
\]
The split controller has underlying Mealy type
\[
\mathbf K_{\mathrm{sgn}}:O\times R\rightsquigarrow (W\times T)+A.
\]
It represents an Int morphism
\[
\mathbf K_{\mathrm{sgn}}:(O\times R,A)\to(W\times T,0).
\]
Their signed feedback composite is
\[
\mathbf K_{\mathrm{sgn}}\circ_{\mathrm{Int}}\mathbf P:(0,0)\to(W\times T,0),
\]
where composition traces over the middle signed interface
\[
(O\times R,A).
\]

\begin{definition}[Abstract signed cybernetic Mealy system]
\label{def:abstract-signed-cybernetic-mealy-system}
An abstract signed cybernetic Mealy system consists of a plant
\[
\mathbf P:A\rightsquigarrow O\times R
\]
and a split adaptive controller
\[
\mathbf K_{\mathrm{sgn}}:O\times R\rightsquigarrow (W\times T)+A,
\]
together with their signed-interface feedback composite
\[
\mathbf K_{\mathrm{sgn}}\circ_{\mathrm{Int}}\mathbf P:(0,0)\to(W\times T,0).
\]
\end{definition}

The trace supplies the feedback wiring.  The stateful Mealy components supply the transition,
output, memory, and adaptation.

\section{Effectful Cybernetic Mealy Systems}
\label{sec:effectful-cybernetic-mealy-systems}

The effectful version replaces the plant and controller by effectful Mealy components.

\begin{definition}[Operational effectful cybernetic Mealy system]
\label{def:operational-effectful-cybernetic-mealy-system}
Let \(\mathsf M\) be a response effect.  An operational \(\mathsf M\)-effectful cybernetic Mealy
system consists of an effectful plant
\[
\mathbf P_{\mathsf M}:A\rightsquigarrow_{\mathsf M}O\times R
\]
and an effectful operational controller
\[
\mathbf K_{\mathsf M,\mathrm{op}}:O\times R\rightsquigarrow_{\mathsf M}A\times(W\times T).
\]
Its decision-update closed-loop semantics is the synchronous effectful process constructed in
Theorem~\ref{thm:closed-loop-effectful-process}.
\end{definition}

\begin{definition}[Signed effectful cybernetic Mealy system]
\label{def:signed-effectful-cybernetic-mealy-system}
Let \(\mathsf M\) be a response effect for which the relevant signed feedback trace is available.
A signed \(\mathsf M\)-effectful cybernetic Mealy system consists of an effectful plant
\[
\mathbf P_{\mathsf M}:A\rightsquigarrow_{\mathsf M}O\times R
\]
and an effectful split controller
\[
\mathbf K_{\mathsf M,\mathrm{sgn}}:O\times R\rightsquigarrow_{\mathsf M}(W\times T)+A,
\]
together with their signed feedback composite
\[
\mathbf K_{\mathsf M,\mathrm{sgn}}\circ_{\mathrm{Int}}\mathbf P_{\mathsf M}.
\]
\end{definition}

The interpretation is:
\[
\boxed{
\text{trace supplies wiring;}
}
\]
\[
\boxed{
\text{the effect supplies one-step uncertainty, branching, failure, cost, or sampling;}
}
\]
\[
\boxed{
\text{Mealy state supplies memory and adaptation.}
}
\]

The deterministic cybernetic systems are the unit image of the effectful theory.  If
\[
\mathsf M=\mathcal D
\]
is the finite distribution effect, then deterministic systems embed as Dirac-valued one-step
systems.

\section{Adaptive Controllers}
\label{sec:adaptive-controllers-effectful}

For the remainder of the concrete cybernetic presentation, we work in the object-level
one-object or discrete-interface specialization.  Thus \(A,O,R,W,T\) denote the action,
observation, reward/performance, value/certificate, and trace objects.  The fully typed
internal-category version replaces these objects by the corresponding object-of-arrows and by the
pullbacks enforcing source-target compatibility.

The adaptive controller is itself stateful.  Its internal state decomposes into controller mode,
policy/model state, optimizer state, and updater memory.

\begin{definition}[Adaptive controller state]
\label{def:adaptive-controller-state-effectful}
An adaptive controller has internal state object
\[
K_0=C\times \Theta\times Z\times U,
\]
where
\[
C=\text{controller mode, guard, or safety state},
\]
\[
\Theta=\text{policy, model, parameter, critic, or belief state},
\]
\[
Z=\text{optimizer state},
\]
and
\[
U=\text{updater memory}.
\]
\end{definition}

The four pieces have different functions.  The controller mode \(C\) records modes, guards, and
safety information.  The policy/model state \(\Theta\) stores the information used to choose
actions.  The optimizer state \(Z\) stores tie-breaking, search state, cache data, gradient state,
or other optimization information.  The updater memory \(U\) stores the information needed to
change future controller behaviour.

\subsection{Guards and admissible actions}
\label{subsec:effectful-guards-admissible-actions}

A controller may restrict the actions that are currently admissible.  This is represented by a
subobject
\[
G\hookrightarrow C\times O\times A.
\]
In finite sets, this is equivalently a family of admissible action sets
\[
G(c,o)\subseteq A.
\]

Let
\[
J=Z\times\Theta\times C\times O.
\]
Pull back the guard along
\[
J\times A\to C\times O\times A
\]
to obtain
\[
G_J\hookrightarrow J\times A.
\]
Thus \(G_J\) is the object of tuples
\[
(z,\theta,c,o,a)
\]
such that \(a\) is admissible in mode \(c\) under observation \(o\).

\subsection{Evaluation}
\label{subsec:effectful-evaluation}

Let \(W\) be a value or certificate object.  An evaluator is a map
\[
Q:\Theta\times C\times O\times A\longrightarrow W.
\]
For fixed \((\theta,c,o)\), the evaluator assigns to each admissible action \(a\) a value
\[
Q(\theta,c,o,a).
\]
The evaluator scores actions; it does not itself have to choose one.

\subsection{Effectful guarded policies}
\label{subsec:effectful-policies}

Let \(T\) be a trace or auxiliary-output object.  Elements of \(T\) may record log-probabilities,
value estimates, gradients, scores, chosen branches, certificates, or other data needed for
learning.

\begin{definition}[Effectful guarded policy]
\label{def:effectful-guarded-policy}
An effectful guarded policy is a response map over \(J\)
\[
\Pi_G:
J\longrightarrow
\mathsf M_J\bigl(G_J\times_J(T\times Z)\bigr).
\]
Elementwise, this is a response
\[
\Pi_G(z,\theta,c,o)
\in
\mathsf M\{(a,t,z^+)\mid a\in G(c,o)\}.
\]
The induced unguarded policy is obtained by composing with the inclusion
\[
G_J\times_J(T\times Z)
\hookrightarrow
J\times A\times T\times Z.
\]
\end{definition}

In deterministic greedy control, the policy is the unit of a selected action, trace, and next
optimizer state:
\[
\Pi(z,\theta,c,o)
=
\eta(a^\ast,t^\ast,z^+).
\]
In probabilistic control, \(\Pi(z,\theta,c,o)\) is a distribution over admissible actions, traces,
and next optimizer states.

\subsection{Optimizers and the max/argmax instance}
\label{subsec:effectful-optimizers}

A deterministic greedy optimizer is a special case of an effectful policy mechanism.  In this
subsection, \(W\) is read in an ordered-value specialization, so that values of
\[
Q:\Theta\times C\times O\times A\to W
\]
can be compared along admissible action fibres.

An optimizer may produce
\[
\Omega:Z\times\Theta\times C\times O
\longrightarrow A\times T\times W\times Z.
\]
Writing
\[
\Omega(z,\theta,c,o)=(a^\ast,t^\ast,w^\ast,z^+),
\]
the intended equations are
\[
a^\ast\in \operatorname{argmax}_{a\in G(c,o)}Q(\theta,c,o,a),
\]
and
\[
w^\ast=\max_{a\in G(c,o)}Q(\theta,c,o,a).
\]

This max/argmax structure is important, but it is not the universal form of optimization.  It is
the deterministic greedy instance.  In stochastic policy-gradient or entropy-regularized control,
one may instead produce a distribution
\[
\Pi(z,\theta,c,o)\in \mathcal D(G(c,o)\times T\times Z)
\]
or update parameters by a gradient-like rule.

The span-valued version avoids choosing a representative maximizer.  The optimizing relation is the
subobject
\[
\operatorname{Opt}_Q
\hookrightarrow
Z\times\Theta\times C\times O\times A\times T\times W\times Z
\]
classifying tuples
\[
(z,\theta,c,o,a,t,w,z^+)
\]
such that \(a\) is admissible, \(w=Q(\theta,c,o,a)\), and \(w\) is maximal among the values
\[
Q(\theta,c,o,b)
\]
for admissible \(b\in G(c,o)\).  A deterministic optimizer is a section of this optimizing
subobject.  Without such a section, the optimizer is naturally span-valued or nondeterministic.

\subsection{Updaters}
\label{subsec:effectful-updaters}

The updater changes the future controller, policy, model, optimizer, and memory state after the
plant consequence is observed.

\begin{definition}[Effectful updater]
\label{def:effectful-updater}
An effectful updater is a response map
\[
\mathcal U:
C\times\Theta\times Z\times U
\times O\times A\times T\times W\times R\times O
\longrightarrow
\mathsf M(C\times\Theta\times Z\times U).
\]
\end{definition}

The input variables are:
\[
c,\theta,z,u
\]
for the current adaptive-controller state,
\[
o
\]
for the observation used to choose the present action,
\[
a
\]
for the selected action,
\[
t
\]
for the policy trace,
\[
w
\]
for the value or certificate,
\[
r
\]
for the resulting reward, cost, or performance signal,
and
\[
o'
\]
for the post-action observation.

The distinction is:
\[
\boxed{
\text{policy/optimizer = within-step decision;}
}
\]
\[
\boxed{
\text{updater = across-step adaptation.}
}
\]
The policy or optimizer chooses using the current policy/model and optimizer state.  The updater
changes that state for future choices.

\section{Closed-Loop Process Presentation}
\label{sec:closed-loop-process-presentation-effectful}

We now derive the autonomous process obtained after closing the plant-controller loop in the
decision-update presentation.

Let the plant state object be \(X\).  Let the adaptive-controller state object be
\[
C\times\Theta\times Z\times U.
\]
The closed-loop state object is
\[
S=X\times C\times\Theta\times Z\times U.
\]

Let
\[
B_{\mathrm{vis}}
\]
be the visible behaviour object.  For reinforcement-learning style systems, a useful choice is
\[
B_{\mathrm{vis}}=O\times A\times R\times W\times T,
\]
recording observation, action, reward or performance signal, value certificate, and policy trace.

\subsection{Deterministic closed loop}
\label{subsec:deterministic-closed-loop-effectful-chapter}

In the deterministic case, suppose the plant has transition, observation, and performance maps
\[
\tau:X\times A\to X,
\]
\[
h:X\to O,
\]
and
\[
\rho:X\times A\times X\to R.
\]
Suppose the controller uses a deterministic optimizer
\[
\Omega:Z\times\Theta\times C\times O\to A\times T\times W\times Z
\]
and deterministic updater
\[
L:
C\times\Theta\times Z\times U
\times O\times A\times T\times W\times R\times O
\to
C\times\Theta\times Z\times U.
\]

Given
\[
s=(x,c,\theta,z,u)\in S,
\]
define
\[
o=h(x).
\]
The optimizer produces
\[
(a^\ast,t^\ast,w^\ast,z^+)=\Omega(z,\theta,c,o).
\]
The plant responds by
\[
x'=\tau(x,a^\ast).
\]
The post-action observation is
\[
o'=h(x').
\]
The performance signal is
\[
r=\rho(x,a^\ast,x').
\]
The updater then produces
\[
(c',\theta',z',u')
=
L(c,\theta,z^+,u,o,a^\ast,t^\ast,w^\ast,r,o').
\]
Thus the closed-loop transition is
\[
F:S\to S,
\]
defined by
\[
F(x,c,\theta,z,u)=(x',c',\theta',z',u').
\]
The visible behaviour map is
\[
b:S\to B_{\mathrm{vis}},
\qquad
b(x,c,\theta,z,u)=(o,a^\ast,r,w^\ast,t^\ast).
\]

\begin{proposition}[Closed deterministic cybernetic process]
\label{prop:closed-deterministic-cybernetic-process}
Every deterministic adaptive cybernetic system of the above form determines an autonomous
deterministic process
\[
F:S\to S
\]
on the coupled state object
\[
S=X\times C\times\Theta\times Z\times U
\]
together with a visible behaviour map
\[
b:S\to B_{\mathrm{vis}}.
\]
\end{proposition}

\begin{proof}
The displayed equations define each component of the next closed-loop state uniquely as a function
of the current closed-loop state.  Hence they define a function \(F:S\to S\).  The same equations
define the visible behaviour map \(b:S\to B_{\mathrm{vis}}\).
\end{proof}

\subsection{Effectful closed loop}
\label{subsec:effectful-closed-loop}

The effectful case replaces deterministic choices by Kleisli composition.

Let the plant be an effectful response
\[
\mathcal P:X\times A\longrightarrow \mathsf M(X\times R).
\]
Let the observation readout be
\[
h:X\to O.
\]
Let the policy be
\[
\Pi:Z\times\Theta\times C\times O
\longrightarrow
\mathsf M(A\times T\times Z).
\]
Let the evaluator be
\[
Q:\Theta\times C\times O\times A\to W.
\]
Let the updater be
\[
\mathcal U:
C\times\Theta\times Z\times U
\times O\times A\times T\times W\times R\times O
\longrightarrow
\mathsf M(C\times\Theta\times Z\times U).
\]

\begin{theorem}[Closed-loop effectful process]
\label{thm:closed-loop-effectful-process}
The above data determine a closed-loop effectful process
\[
K:S\to \mathsf M(S\times B_{\mathrm{vis}})
\]
on the coupled state object
\[
S=X\times C\times\Theta\times Z\times U,
\]
where
\[
B_{\mathrm{vis}}=O\times A\times R\times W\times T.
\]
\end{theorem}

\begin{proof}
The closed-loop process is the Kleisli composite
\[
K:S\to \mathsf M(S\times B_{\mathrm{vis}})
\]
defined by the following do-notation:
\[
\begin{aligned}
K(x,c,\theta,z,u)
=
{}&
\text{\rm let } o=h(x) \text{\rm ;}\\
&
(a,t,z^+)\leftarrow \Pi(z,\theta,c,o);\\
&
\text{\rm let } w=Q(\theta,c,o,a) \text{\rm ;}\\
&
(x',r)\leftarrow \mathcal P(x,a);\\
&
\text{\rm let } o'=h(x') \text{\rm ;}\\
&
(c',\theta',z',u')
\leftarrow
\mathcal U(c,\theta,z^+,u,o,a,t,w,r,o');\\
&
\eta\bigl((x',c',\theta',z',u'),(o,a,r,w,t)\bigr).
\end{aligned}
\]
The policy, plant, and updater are composable Kleisli arrows for \(\mathsf M\).  The observation
readout and evaluator are ordinary deterministic maps and therefore enter through the unit of the
effect.  By associativity of Kleisli composition, the displayed composite is a well-defined map
\[
S\to \mathsf M(S\times B_{\mathrm{vis}}).
\]
\end{proof}

\begin{corollary}[Unit embedding of deterministic closed loops]
\label{cor:unit-embedding-deterministic-closed-loops}
The deterministic closed-loop transition and visible behaviour
\[
F:S\to S,
\qquad
b:S\to B_{\mathrm{vis}}
\]
embed into the effectful theory by
\[
K_F(s)=\eta(F(s),b(s)).
\]
\end{corollary}

\begin{proof}
This is the image of the deterministic closed-loop process under the unit of \(\mathsf M\).
\end{proof}

\begin{remark}[Observation-emitting plants]
\label{rem:observation-emitting-plants}
The presentation above uses a deterministic readout
\[
h:X\to O
\]
and an effectful plant response
\[
\mathcal P:X\times A\to\mathsf M(X\times R).
\]
A partially observed or sensor-noisy variant replaces this by an observation-emitting plant kernel
\[
\mathcal P:X\times A\to\mathsf M(X\times O\times R).
\]
The same Kleisli construction applies, with the sampled observation replacing \(h(x')\).
\end{remark}

\section{Trace-Compatible Effectful Feedback}
\label{sec:trace-compatible-effectful-feedback}

The operational closed loop above is a synchronous one-cycle feedback construction.  Chapter~9's
feedback trace is a different but related construction: it hides a coproduct summand by finite
hidden-path closure.

For spans, a morphism
\[
R:X+H\to Y+H
\]
has block decomposition
\[
R=
\begin{pmatrix}
R_{XY} & R_{XH}\\
R_{HY} & R_{HH}
\end{pmatrix},
\]
and the coproduct trace is
\[
\Tr_H(R)
=
R_{XY}
\boxplus
R_{HY}\circ \Path(R_{HH})\circ R_{XH}.
\]
The hidden interface \(H\) is traversed for finitely many hidden steps before exiting.

In an effectful setting, hidden paths must be composed using Kleisli composition.  This is not
automatic for every response effect.

\begin{definition}[Trace-compatible response effect]
\label{def:trace-compatible-response-effect}
A response effect \(\mathsf M\) is trace-compatible on a class of effectful spans if the following
operations are available inside that class.

\begin{enumerate}
\item Kleisli composition of effectful spans, denoted
\[
\odot .
\]

\item Finite Kleisli powers of every hidden endospan
\[
U:H\to H,
\qquad
U^{\odot n}:H\to H.
\]

\item A finite-path object
\[
\Path_{\mathsf M}(U)
=
\coprod_{n\geq 0}U^{\odot n}
\]
representing finite hidden Kleisli executions.

\item For every block effectful span
\[
R:X+H\to Y+H
\]
with block decomposition
\[
R=
\begin{pmatrix}
R_{XY} & R_{XH}\\
R_{HY} & R_{HH}
\end{pmatrix},
\]
the expression
\[
\Tr_H^{\mathsf M}(R)
=
R_{XY}
\boxplus
R_{HY}\odot
\Path_{\mathsf M}(R_{HH})
\odot
R_{XH}
\]
is again an effectful span of the same class.
\end{enumerate}
\end{definition}

\begin{proposition}[Effectful hidden-path trace]
\label{prop:effectful-hidden-path-trace}
For a trace-compatible response effect \(\mathsf M\), every effectful span
\[
R:X+H\to Y+H
\]
in the trace-compatible class determines a traced effectful span
\[
\Tr_H^{\mathsf M}(R):X\to Y
\]
by finite Kleisli hidden-path closure.
\end{proposition}

\begin{proof}
The trace-compatible structure supplies Kleisli powers
\[
R_{HH}^{\odot n}
\]
and their finite-path closure
\[
\Path_{\mathsf M}(R_{HH})=\coprod_{n\geq 0}R_{HH}^{\odot n}.
\]
Composing first from \(X\) into \(H\), then along a finite hidden path, and finally from \(H\) to
\(Y\), and adding the direct \(X\to Y\) component, gives the displayed formula for
\(\Tr_H^{\mathsf M}(R)\).
\end{proof}

For nondeterminism, this collects all finite hidden behaviours.  For subprobability, it sums the
mass of finite hidden trajectories that eventually exit.  For the finite-distribution monad,
unbounded hidden-path trace is not automatically a finite distribution.  It is available in guarded
finite-horizon, nilpotent, or finitely exiting cases, and also after moving to a countable or
subprobabilistic kernel setting in which the total mass of finite exiting trajectories is defined.
Otherwise the probabilistic operational closure should be read as a one-step Markov kernel rather
than as an unbounded coproduct trace.

Thus, for probabilistic cybernetics, it is important to distinguish:
\[
\text{synchronous one-cycle closed interaction}
\]
from
\[
\text{episode-level guarded trace}.
\]
The first gives ordinary Markov-kernel dynamics.  The second gives finite rollout or guarded-while
semantics.

\section{Probabilistic Cybernetic Mealy Systems}
\label{sec:probabilistic-cybernetic-mealy-systems}

We now instantiate the response effect by finite probability.  Let
\[
\mathcal D(X)
\]
denote the set of finite probability distributions on \(X\).  In this case an effectful Mealy
process has response
\[
D_\Phi\to \mathcal D_{D_\Phi}(N_\Phi),
\]
or elementwise
\[
\kappa_\Phi(a,x)
\in
\mathcal D
\{(y,\beta):p(y)=s_A(a),\;\beta:q(y)\to q(x)\}.
\]
The deterministic Mealy morphism embeds by Dirac distributions:
\[
(a,x)\mapsto
\delta_{(\delta_\Phi(a,x),\omega_\Phi(a,x))}.
\]

\begin{definition}[Finite probabilistic adaptive cybernetic system]
\label{def:finite-probabilistic-adaptive-cybernetic-system}
A finite probabilistic adaptive cybernetic system consists of:
\[
X,\ C,\ \Theta,\ Z,\ U
\]
for plant, controller, policy/model, optimizer, and updater state;
interfaces
\[
A,\ O,\ R,\ W,\ T;
\]
an observation readout
\[
h:X\to O;
\]
a guard relation
\[
G\hookrightarrow C\times O\times A;
\]
a plant kernel
\[
\mathcal P:X\times A\to \mathcal D(X\times R);
\]
a policy kernel
\[
\Pi:Z\times\Theta\times C\times O\to \mathcal D(A\times T\times Z)
\]
whose action component is supported on admissible actions;
an evaluator
\[
Q:\Theta\times C\times O\times A\to W;
\]
and an updater kernel
\[
\mathcal U:
C\times\Theta\times Z\times U
\times O\times A\times T\times W\times R\times O
\to
\mathcal D(C\times\Theta\times Z\times U).
\]
\end{definition}

The closed-loop state is
\[
S=X\times C\times\Theta\times Z\times U.
\]
The visible behaviour object is usually taken to be
\[
B_{\mathrm{vis}}=O\times A\times R\times W\times T.
\]

\begin{theorem}[Closed-loop Markov kernel]
\label{thm:closed-loop-markov-kernel}
Every finite probabilistic adaptive cybernetic system determines a Markov kernel
\[
K:S\to \mathcal D(S\times B_{\mathrm{vis}})
\]
on the coupled state object
\[
S=X\times C\times\Theta\times Z\times U.
\]
If the plant, policy, and updater kernels are Dirac-valued, then \(K\) is Dirac-valued and the
system is deterministic.
\end{theorem}

\begin{proof}
The kernel is the Kleisli composite for the finite distribution monad:
\[
\begin{aligned}
K(x,c,\theta,z,u)
=
{}&
\text{\rm let } o=h(x) \text{\rm ;}\\
&
(a,t,z^+)\sim \Pi(z,\theta,c,o);\\
&
\text{\rm let } w=Q(\theta,c,o,a) \text{\rm ;}\\
&
(x',r)\sim \mathcal P(x,a);\\
&
\text{\rm let } o'=h(x') \text{\rm ;}\\
&
(c',\theta',z',u')
\sim
\mathcal U(c,\theta,z^+,u,o,a,t,w,r,o');\\
&
\delta_{\bigl((x',c',\theta',z',u'),(o,a,r,w,t)\bigr)} .
\end{aligned}
\]
Since \(\mathcal D\) is a monad, sequencing the policy, plant, and updater kernels gives a
probability distribution on the resulting closed-loop state and visible behaviour.  If every
component kernel is Dirac-valued, the composite distribution is Dirac-valued.
\end{proof}

This is the Markov-process object canonically produced by the probabilistic cybernetic pattern.
Before feedback closure, the system is an open stochastic controlled process.  After operational
closure, it is an autonomous Markov kernel on the coupled state object.

\subsection{Relation to controlled Markov processes}
\label{subsec:relation-to-mdps}

The probabilistic Mealy process has the shape of a typed stochastic controlled process:
\[
\text{state}+\text{admissible typed action}
\longmapsto
\text{distribution over next state and structured output}.
\]
In the one-object finite-set case this resembles the usual controlled kernel
\[
X\times A\to \mathcal D(X\times R)
\]
together with an observation readout
\[
h:X\to O.
\]
The general typed version remembers more structure:
\[
D_\Phi=A_1\times_{t_A,A_0,p}P
\]
is the admissible state-action object, and the outputs are arrows in an output interface rather
than unstructured labels.

A Markov reward process is obtained only after the reward or performance interface \(R\) is chosen
as part of the output, or after a further readout map
\[
B_1\to R
\]
is supplied.  Thus reward is not inherent in the Mealy structure; it is an interpretation of part
of the structured output.

\section{Examples}
\label{sec:effectful-cybernetic-examples}

\subsection{Dirac deterministic thermostat}
\label{subsec:dirac-deterministic-thermostat}

This example illustrates the deterministic unit embedding.

Let the plant state be
\[
X=\{0,1,2,3,4\}.
\]
Let the action set be
\[
A=\{H,I,C\},
\]
where \(H\) heats, \(I\) idles, and \(C\) cools.  Define
\[
\tau(x,H)=\min(4,x+1),
\]
\[
\tau(x,I)=x,
\]
and
\[
\tau(x,C)=\max(0,x-1).
\]
Let the observation be full-state observation:
\[
h(x)=x.
\]
Let the performance signal be
\[
\rho(x,a,x')=-|x'-2|.
\]

A deterministic plant kernel is the Dirac kernel
\[
\mathcal P(x,a)
=
\delta_{(\tau(x,a),\rho(x,a,\tau(x,a)))}.
\]
The post-action observation is then recovered as
\[
h(\tau(x,a)).
\]
A deterministic greedy policy is represented by a Dirac policy kernel
\[
\Pi(z,\theta,c,o)=\delta_{(a^\ast,t^\ast,z^+)}
\]
where \(a^\ast\) maximizes a chosen evaluator
\[
Q(\theta,c,o,a).
\]
The closed-loop stochastic process is therefore Dirac-valued.

\subsection{Stochastic thermostat}
\label{subsec:stochastic-thermostat}

This example illustrates a probabilistic plant kernel.

A stochastic variant keeps the same action set but allows plant noise.  After choosing the intended
next temperature
\[
\bar x=\tau(x,a),
\]
the realized next temperature may be
\[
x'\in\{\max(0,\bar x-1),\bar x,\min(4,\bar x+1)\}
\]
with probabilities depending on external disturbance.  Thus the plant is now a genuine kernel
\[
\mathcal P:X\times A\to \mathcal D(X\times R),
\]
where the reward component is computed from the realized \(x'\).  The policy may remain
deterministic, or it may sample among admissible actions.  The feedback-closed system is then a
Markov kernel on
\[
S=X\times C\times\Theta\times Z\times U.
\]

\subsection{Stochastic bandit learning}
\label{subsec:stochastic-bandit-learning}

This example illustrates terminal observation and adaptive updating.

Let
\[
A=\{0,1\}
\]
be the action set.  The observation object is terminal:
\[
O=1.
\]
Let the reward object be
\[
R=\{0,1\}.
\]
A stochastic bandit plant is a kernel
\[
\mathcal P:X\times A\to \mathcal D(X\times R).
\]
For a stationary two-armed bandit, the plant state \(X\) may be terminal and the kernel is
determined by reward probabilities
\[
\Pr(r=1\mid 0)=\mu_0,
\qquad
\Pr(r=1\mid 1)=\mu_1.
\]

Let the policy/model state \(\Theta\) store estimated values
\[
\theta=(\hat r_0,\hat r_1).
\]
A greedy deterministic policy chooses
\[
i^\ast\in\operatorname{argmax}_{i\in A}\hat r_i.
\]
An exploratory stochastic policy may instead use
\[
\Pi(z,\theta,c,\ast)\in\mathcal D(A\times T\times Z),
\]
for example an \(\epsilon\)-greedy or softmax distribution.

The updater changes the estimate of the selected arm after observing the reward:
\[
\theta'\sim \mathcal U(c,\theta,z,u,\ast,i,t,w,r,\ast).
\]
The deterministic table update is the Dirac special case.  Stochastic approximation, randomized
exploration schedules, and noisy update rules are represented by non-Dirac updater kernels.

\subsection{Inventory regulation with stochastic demand}
\label{subsec:stochastic-inventory-regulation}

This example illustrates a stochastic plant with model update.

Let inventory levels be
\[
I=\{0,1,2,3,4,5\}
\]
and actions be order quantities
\[
A=\{0,1,2,3\}.
\]
The plant state is the current inventory level \(x\).  Let the true demand be sampled from a
plant-side distribution
\[
d\sim \nu_x,
\]
or, in the stationary case,
\[
d\sim \nu.
\]
Given order \(a\), the stocked amount is
\[
y=\min(5,x+a).
\]
The next inventory is
\[
x'=\max(0,y-d).
\]
The reward may be negative cost:
\[
r=-(2\max(0,d-y)+\max(0,y-d)+a).
\]
Thus
\[
\mathcal P(x,a)
\in
\mathcal D(I\times R)
\]
is induced by the demand distribution.  The policy state \(\Theta\) may store an estimate of the
demand law, the controller state \(C\) may store cooldown or budget mode, and the updater changes
\(\Theta\) after observing realized demand.

\subsection{Concurrent plants under a shared adaptive controller}
\label{subsec:concurrent-plants-effectful}

This example illustrates product tensor together with a coupled controller guard.

Let two plants have state objects \(X_1,X_2\), action interfaces \(A_1,A_2\), and plant kernels
\[
\mathcal P_i:X_i\times A_i\to \mathsf M(X_i\times R_i).
\]
Their parallel plant has state object
\[
X=X_1\times X_2
\]
and action object
\[
A=A_1\times A_2.
\]

For a general response effect, forming the product plant requires a product or Fubini operation
\[
\mathsf M(X)\times\mathsf M(Y)\to\mathsf M(X\times Y).
\]
For independent finite probability, this is the ordinary product distribution:
\[
\mathcal P_1\otimes \mathcal P_2.
\]

A shared controller may impose a coupled guard
\[
G\hookrightarrow C\times O_1\times O_2\times A_1\times A_2,
\]
such as an energy or resource constraint.  The plant dynamics remain parallel; the controller
introduces coupled decision-making through the guard and evaluator.  Thus causal dependence through
the controller does not collapse the plant-side product into interleaving semantics.

\section{Specifications and Behaviour}
\label{sec:effectful-cybernetic-specifications}

After operational feedback closure, an effectful cybernetic system gives
\[
K:S\to \mathsf M(S\times B_{\mathrm{vis}}).
\]
State specifications are subobjects
\[
\varphi\hookrightarrow S.
\]
For relational or nondeterministic effects, the diamond modality expresses possible reachability
and the box modality expresses necessary preservation.  For probabilistic effects, one may instead
interpret specifications by probability thresholds, expectation, almost-sure properties, or
subprobability mass.

In the deterministic unit case
\[
K(s)=\eta(F(s),b(s)),
\]
the one-step diamond and box reduce to inverse image along \(F\).  Let
\[
U_F:S\to S
\]
be the endospan associated to \(F:S\to S\).  Finite-execution reachability and invariance are
expressed by
\[
\langle \Path(U_F)\rangle\varphi
\]
and
\[
[\Path(U_F)]\varphi.
\]
Thus one-step modalities are pullback modalities along \(F\), while path modalities are finite
iteration modalities.

Visible behaviour is obtained by projecting to
\[
B_{\mathrm{vis}}.
\]
In the deterministic case, an initial state \(s_0\) produces a stream
\[
b(s_0),\ b(F(s_0)),\ b(F^2(s_0)),\ldots .
\]
In the probabilistic case, an initial state induces a distribution over visible behaviour streams,
or over finite traces when a guarded finite horizon is chosen.

\section{Chapter Summary}
\label{sec:effectful-cybernetic-summary}

The constructions of this chapter form the following hierarchy.

\begin{enumerate}
\item Every deterministic Mealy morphism
\[
\Phi:A\rightsquigarrow B
\]
has a typed one-step process presentation
\[
D_\Phi\to N_\Phi,
\]
where
\[
D_\Phi=A_1\times_{t_A,A_0,p}P
\]
is the admissible state-input object and \(N_\Phi\to D_\Phi\) is the typed next-output object.

\item A one-step process presentation is a Mealy morphism precisely when it satisfies the identity
and composite-input Mealy laws:
\[
\delta(1_{p(x)},x)=x,
\qquad
\omega(1_{p(x)},x)=1_{q(x)},
\]
and
\[
\delta(a\circ b,x)
=
\delta\bigl(b,\delta(a,x)\bigr),
\]
\[
\omega(a\circ b,x)
=
\omega(a,x)\circ
\omega\bigl(b,\delta(a,x)\bigr).
\]

\item A response-effect context is an indexed Kleisli doctrine sending a typed response family
\[
N\to D
\]
to an effectful response family
\[
\mathsf M_D(N)\to D
\]
and supporting dependent Kleisli substitution.

\item An effectful Mealy morphism
\[
\Phi:A\rightsquigarrow_{\mathsf M}B
\]
is an effectful one-step response
\[
D_\Phi\to\mathsf M_{D_\Phi}(N_\Phi)
\]
satisfying the Kleisli identity law and the dependent Kleisli composite-input law over the
composable state-input object
\[
C_\Phi
=
\bigl(A_1\times_{s_A,A_0,t_A}A_1\bigr)
\times_{t_A\pi_1,A_0,p}
P .
\]

\item Deterministic Mealy morphisms embed into effectful Mealy morphisms by the unit of the
response effect:
\[
(a,x)\mapsto
\eta(\delta_\Phi(a,x),\omega_\Phi(a,x)).
\]

\item When an effectful output-comparison Kleisli structure is available, including a lift of
\[
R_A=-\circ A
\]
to
\[
\Kl(T_B^{\mathsf M}),
\]
effectful Mealy morphisms are \(R_A\)-algebras in
\[
\Kl(T_B^{\mathsf M}).
\]

\item Strict effectful simulations are oppositely directed carrier maps commuting with effectful
response maps.  A strict simulation
\[
\Theta:\Phi\Rightarrow\Psi
\]
with carriers \(P\) and \(Q\) has carrier component
\[
\theta:Q\to P.
\]
More general output-lax simulations are obtained as Kleisli algebra morphisms in the same
effectful output-comparison setting.

\item Operational cybernetic feedback produces a closed-loop effectful process
\[
K:S\to\mathsf M(S\times B_{\mathrm{vis}})
\]
on the coupled state object
\[
S=X\times C\times\Theta\times Z\times U.
\]

\item Signed-interface cybernetic feedback uses the coproduct trace.  The controller must be
presented in split form
\[
O\times R\rightsquigarrow (W\times T)+A
\]
to represent an Int morphism
\[
(O\times R,A)\to(W\times T,0).
\]

\item Trace feedback is available for trace-compatible response effects and is computed by
effectful finite hidden-path closure:
\[
\Tr^{\mathsf M}_H(R)
=
R_{XY}
\boxplus
R_{HY}\odot
\Path_{\mathsf M}(R_{HH})
\odot
R_{XH}.
\]

\item In the finite probabilistic specialization, the operational closed loop is a Markov kernel
\[
K:S\to\mathcal D(S\times B_{\mathrm{vis}}).
\]
If the plant, policy, and updater kernels are Dirac-valued, this Markov kernel is deterministic.
\end{enumerate}

Thus probabilistic reinforcement-learning style interaction is obtained not by replacing the Mealy
theory, but by choosing the finite distribution response effect and closing the resulting effectful
Mealy network by cybernetic feedback.